\section{Prasanth et al. (2023)}

\subsection*{Resumo}

O artigo discute como a inteligência artificial é usada para apoiar a tomada de decisão nas empresas e como ela vem alterando modelos de negócio (p.~965). Segundo os autores, sistemas de IA permitem processar grandes volumes de dados e tomar decisões mais rápidas e mais bem informadas, com ganhos de eficiência, precisão e inovação (p.~965).

A revisão de literatura relaciona a IA a \emph{big data}, \emph{business intelligence} e \emph{business analytics}, e cita aplicações em comércio eletrônico, operações de \emph{trading}, detecção de fraude e análise de informações financeiras (p.~966). Os autores também apresentam uma estrutura conceitual com três formas de combinar humanos e IA na decisão: decisão agregada, delegação total à IA e decisão sequencial, de humano para IA ou de IA para humano (p.~967).

Nos resultados, os autores afirmam que a IA pode apoiar, substituir ou aprimorar a decisão humana na formulação de planos de marketing, e que a modelagem preditiva permite simular como uma ação afetaria diferentes partes da organização (p.~968). Ao mesmo tempo, apontam preocupações com privacidade, segurança, transparência, ética e impacto sobre o trabalho (p.~965 e~968).

A conclusão afirma que a IA permite decisões melhores ao analisar grandes volumes de dados e gerar previsões e sugestões, desde que usada de forma responsável e transparente  (p.~968).

\subsection*{Natureza do texto}

Texto conceitual, com revisão de literatura narrativa. A seção chamada ``Methodology'' (p.~967) não descreve dados nem procedimentos, e as figuras e a tabela apresentadas como resultados (p.~968) não indicam a origem dos dados. Há também inconsistências: os resultados dizem que a IA pode substituir a decisão humana, mas a conclusão nega essa ideia (p.~968), e algumas citações numeradas não correspondem à lista de referências (p.~966--969). Serve como contexto conceitual sobre IA e decisão.

\subsection*{O que o texto trouxe de novo}

O principal aprendizado foi entender melhor a diferença entre previsão e decisão. Um modelo pode prever uma variável ou identificar um padrão, mas isso não significa que a decisão já esteja determinada. Ainda existe uma etapa em que a informação produzida pelo modelo precisa ser usada para escolher uma ação, e essa etapa depende de algum objetivo, restrição ou consequência.

Isso me ajudou a separar duas abordagens. A modelagem preditiva tenta responder ``o que pode acontecer?''; a modelagem prescritiva usa a informação disponível para responder ``o que pode ser feito?''. Os autores não usam esses termos, mas a distinção aparece quando afirmam que a IA gera previsões e também sugestões (p.~968), e quando descrevem formas diferentes de dividir a decisão entre humanos e IA (p.~967).

Também ficou mais claro que avaliar uma aplicação de IA só pela precisão pode ser insuficiente quando o resultado do modelo interfere em decisões reais.

\subsection*{O que extrair para um trabalho autoral}

O texto pode ser usado como uma das bases conceituais para conectar as duas partes do tema do projeto: a modelagem preditiva e a prescritiva. A partir dele, a pergunta científica deixa de ser só ``o modelo consegue prever?'' e passa a incluir ``a informação produzida pelo modelo melhora uma decisão?''. Essa formulação é minha, e não dos autores.

Também é útil a ideia de Vidgen, Shaw e Grant, citada pelos autores, de que a análise de dados funciona como mediadora entre os dados da organização e o valor econômico obtido com melhores ações e decisões (p.~966). Antes de usar essa ideia no projeto, é preciso ler o texto original.

Por fim, os autores lembram que a decisão com IA envolve custos, riscos e transparência (p.~965 e~967--968). No projeto, isso aparece na forma decustos de transação, controle de risco e registro de todas as tentativas.

\subsection*{Autores para catalogar}

Citados pelos autores, ainda não lidos (status C). Como parte da numeração das citações não confere com a lista de referências, a busca deve ser feita por nome e ano.

\begin{itemize}
  \item Trunk, Birkel e Hartmann (2020) --
        combinação de inteligência humana e
        artificial em decisões estratégicas;
        relevância alta.
  \item Duan, Edwards e Dwivedi (2019) -- IA e
        decisão; relevância alta (também citado em
        Bharadiya et~al.).
  \item Vidgen, Shaw e Grant (2017) -- valor a partir
        de \emph{business analytics}; relevância
        média.
  \item Cavalcante et~al.~(2016) -- inteligência
        computacional em mercados financeiros;
        relevância média.
  \item Xing, Cambria e Welsch (2018) -- previsão
        financeira com linguagem natural; relevância
        baixa por enquanto (extensão P4).
\end{itemize}